% Principal GDZ PPN236018647,native480–484upper/printed473–477upper;HKUST481–485.
\clearpage\SourceLeaf{480}\NoSourceHead
\pdfbookmark[2]{Bestimmung des kleinsten Werthes der Summe}{manifolds-II}
\begin{center}[II.]\end{center}
\HistoricTitle{BESTIMMUNG DES KLEINSTEN WERTHES DER SUMME}
\[
x_1^2+x_2^2+\mathinner{\cdotp\cdotp\cdotp}+x_n^2=R^2
\]
\HistoricTitle{FÜR \(m\) GEGEBENE UNGLEICHUNGEN \(u\geqq 0\).}
\HistoricCitation{[Aus der Vorlesung \textit{Über die Methode der kleinsten Quadrate}, Wintersemester 1850/51, nachgeschrieben von \textsc{August Ritter}.]}
\begin{center}[1.]\end{center}
\indent Die gegebenen Ungleichungen seien linear und ihre Anzahl grösser als die der Variabeln. Für diesen Fall ist das Verfahren folgendes. Wir setzen \(n\) von den Ungleichungen gleich 0 und zwar eine solche Combination von \(n\) Ungleichungen, bei welcher die übrigen \(m-n\) positiv werden, d. h. wir suchen für die Grössen
\[
x_1,\ x_2,\ .\ .\ .,\ x_n
\]
\noindent ein System von Werthen
\[
k_1,\ k_2,\ .\ .\ .,\ k_n,
\]
\noindent welche \(n\) von den Ungleichungen zu 0, die übrigen \(m-n\) positiv werden lassen. Von diesem bestimmten Werthen-System gehen wir nun aus und untersuchen zunächst, wie man dasselbe ändern muss, damit die Summe ihrer Quadrate am schnellsten abnehme, oder, was dasselbe bedeutet, wie man eine vorgeschriebene Abnahme jener Summe durch eine möglichst geringe Aenderung des Werthen-Systems erreichen kann. Als Maass der letzteren betrachten wir die Wurzel aus der Summe der Quadrate der einzelnen Aenderungen. Unter allen Combinationen \(\delta x_1,\ \delta x_2,\ .\ .\ .,\ \delta x_n\), welche eine vorgeschriebene
\par\vfill\noindent X1.\hfill60
\clearpage\SourceLeaf{481}\ArithmeticHead{474}{L}{GEOMETRIE. NACHLASS.}
\noindent Aenderung \(-2\lambda_1\) unserer Funktion bewirken, suchen wir also diejenige heraus, bei welcher die Grösse
\[
\sqrt{\delta x_1^2+\mathinner{\cdotp\cdotp\cdotp}+\delta x_n^2}
\]
\noindent den kleinsten Werth annimmt. Um dies analytisch auszudrücken, haben wir die Variation der Grösse \(R^2\) gleich \(-2\lambda_1\) zu setzen, nachdem wir darin die Werthe \(k_1,.\ .\ .,k_n\) für die Grössen \(x_1,.\ .\ .,x_n\) substituirt haben, und dann zu untersuchen, wie unter dieser Bedingung und unter den aus der Variation von
\[
u_1=0,\quad u_2=0,\quad .\ .\ .,\quad u_n=0
\]
\noindent hervorgehenden Bedingungen
\[
\delta u_1\geqq 0,\quad\delta u_2\geqq 0,\quad .\ .\ .,\quad\delta u_n\geqq 0
\]
\noindent die Grösse
\[
\delta x_1^2+\delta x_2^2+\mathinner{\cdotp\cdotp\cdotp}+\delta x_n^2=R_1^2
\]
\noindent ein Minimum wird. Wir haben also jetzt in Bezug auf die Grössen \(\delta x_1,.\ .\ .,\delta x_n\) genau dieselbe Aufgabe zu lösen wie vorher in Bezug auf \(x_1,.\ .\ .,x_n\), nur dass wir statt wie vorhin \(n\) jetzt nur \(n-1\) Variabeln haben (denn Eine kann mittelst der Gleichung
\[
\delta(R^2)=-2\lambda_1
\]
\noindent eliminirt werden) und statt der \(m\) früheren Ungleichungen nur \(n\). Wenden wir daher dasselbe Verfahren auf diese neue Aufgabe an. —

\indent Wir suchen wieder zuerst eine Combination von Werthen für \(\delta x_1,.\ .\ .,\delta x_n\), welche \(n-1\) von den Ungleichungen zu 0 und die letzte positiv macht und die Gleichung
\[
\delta(R^2)+2\lambda_1=0
\]
\noindent [befriedigt]; seien diese Werthe
\[
l_1,\ l_2,\ .\ .\ .,\ l_n.
\]
\noindent Von dieser Combination von Werthen ausgehend suchen wir dann die kleinste Aenderung dieses Werthen-Systems, welche eine vorgeschriebene Abnahme der Funktion \(R_1^2\) um \(2\lambda_2\) hervorbringt. Die Aenderungen \([\xi_1,\xi_2,.\ .\ .,\xi_n]\) der einzelnen Werthe müssen dabei also folgenden Bedingungen genügen:
\[
\delta^2(R^2)=0,\quad\delta(R_1^2)=-2\lambda_2;\quad\delta^2u_1\geqq0,\quad\delta^2u_2\geqq0,\quad .\ .\ .,\quad\delta^2u_{n-1}\geqq0,
\]
\clearpage\SourceLeaf{482}\ArithmeticHead{475}{R}{MANNIGFALTIGKEITEN VON \(n\) DIMENSIONEN.}
\noindent und wie schon bemerkt soll die Summe ihrer Quadrate \(\xi_1^2+\xi_2^2+\mathinner{\cdotp\cdotp\cdotp}+\xi_n^2=R_2^2\) ein Minimum dabei sein. Dies ist wiederum eine der ursprünglichen völlig analoge Aufgabe; wir haben aber statt der \(n\) Variabeln jetzt nur \(n-2\) (weil zwei von ihnen aus den beiden Gleichungen \([\delta R^2+2\lambda_1=0,\ \delta R_1^2+2\lambda_2=0]\) eliminirt werden können) und statt der \(m\) Ungleichungen nur \(n-1\).

\indent Die Auflösung dieser dritten Aufgabe führt wieder auf eine ähnliche vierte, diese auf eine fünfte etc., wobei jedesmal die folgende Eine Veränderliche und Eine Ungleichung weniger enthält als die ihr vorhergehende Aufgabe.

\indent Im weiteren Verlauf dieses Verfahrens können nun zwei verschiedene Fälle eintreten: entweder es lässt sich bis zum letzten Schritte, wo man es nur noch mit Einer Veränderlichen und 2 Ungleichungen zu thun hat, ohne Hinderniss verfolgen; oder aber es tritt schon vorher ein solches ein, welches der Fortsetzung des Verfahrens in der bisherigen Weise ein Ende macht.

\indent Bei jeder neuen Aufgabe tritt nämlich, wie wir gesehen haben, die Frage auf: wie man durch eine möglichst geringe Aenderung des vorliegenden Werthen-Systems eine gewisse Abnahme der Funktion \(R^2\) um \(2\lambda\) erreichen könne. Es kann nun der Fall eintreten — und darin besteht eben das erwähnte mögliche Hinderniss —, dass eine solche Abnahme der Grösse \(R^2\) überhaupt durch keine Aenderung des Werthen-Systems erreicht werden kann, eben weil sie durch Substitution der betreffenden Werthe schon ihren kleinsten Werth erhalten hat.

\indent In beiden Fällen ist man zu gewissen End-Werthen gelangt, welche der letzten Partial-Aufgabe genügen, und es fragt sich nun, wie man diese Werthe zur Auflösung der Haupt-Aufgabe weiter benutzen kann. Zuvor ist jedoch noch zu bemerken, dass in dem ersten Falle, wo die Natur der Aufgabe es mit sich bringt, dass man das angegebene Verfahren bis zu dem letzten Schritte verfolgen muss, eben in Bezug auf diesen letzten Schritt das Verfahren — wenn auch an sich dasselbe bleibend — doch in seiner Darstellung etwas anders aussehen wird. Bevor man nämlich zur Lösung der letzten Aufgabe schreitet, wo man also \(n-1\) Gleichungen und 2 Ungleichungen hat für die Grössen, deren Minimalwerthe man sucht, hat man zu erwägen, dass nunmehr die Aenderung des Werthen-Systems nur noch in einer einzigen Weise ge-
\par\vfill\noindent\hfill60\textsuperscript{*}
\clearpage\SourceLeaf{483}\ArithmeticHead{476}{L}{GEOMETRIE. NACHLASS.}
\noindent schehen kann, da durch die Aenderung Eines Werthes die aller übrigen schon bedingt ist. Die Frage nach der Art und Weise der Aenderung des Werthen-Systems, wie sie bei allen vorhergehenden Schritten auftrat, ist hier also überflüssig, und es handelt sich hier nur um die Grösse dieser Aenderung, und diese ergiebt sich daraus, dass das zugehörige letzte \(\lambda\), welches eine Funktion derselben ist, Null werden muss, wenn nicht schon vorher die Aenderung den Punkt erreicht, wo die letzte übriggebliebene Ungleichung aufhört erfüllt zu werden. In beiden Fällen erhält \(\lambda\) einen bestimmten Werth, und die Werthe, welche die vorhergehende Minimal-Gleichung erfüllen, sind demnach gefunden.

\indent Die Benutzung dieser Werthe zur Lösung der Haupt-Aufgabe geschieht nun in folgender Weise. Das Erfülltsein der letzten Minimal-Gleichung hat für die ihr vorhergehende die Bedeutung: dass sie die Art und Weise angiebt, wie man die in der vorhergehenden substituirten bestimmten Werthe verändern muss, um sie am schnellsten abnehmen zu lassen. Um sie nun zu einem wirklichen Minimum zu machen, brauchen wir nur die Werthe wirklich in der angegebenen Weise so weit zu ändern, bis entweder die entsprechende Grösse \(\lambda\) zu 0 wird oder die überzählige positiv gelassene Ungleichheit aufhört erfüllt zu sein. Die hierdurch erhaltenen Werthe geben dann an, in welcher Weise das in der vorhergehenden Minimal-Gleichung substituirte Werthen-System geändert werden muss, um sie am schnellsten abnehmen zu lassen, wodurch man in den Stand gesetzt wird, auch sie zu einem wirklichen Minimum zu machen; und so kann man rückwärtsgehend Eine Minimal-Gleichung nach der andern zu einem wirklichen oder relativen Minimum machen, bis man endlich auch die Haupt-Gleichung
\[
x_1^2+x_2^2+\mathinner{\cdotp\cdotp\cdotp}+x_n^2
\]
\noindent zu einem solchen gemacht hat.

\indent Die Aufgabe ist gelöst, wenn bei zunehmender Aenderung der Werthe \(x_1,.\ .\ .,x_n\) die Variation von \(R^2\) Null wird; sie ist dagegen noch nicht erledigt, wenn die Aenderung des Werthen-Systems durch das Nullwerden Einer oder mehrerer der vorhin positiv gelassenen \(m-n\) Ungleichungen aufgehalten wird. In diesem Falle ist das ganze bis dahin beschriebene Verfahren noch einmal in Bezug auf die jetzt Null gewordenen Ungleichungen zu wiederholen.
\clearpage\SourceLeaf{484}\ArithmeticHead{477}{R}{MANNIGFALTIGKEITEN VON \(n\) DIMENSIONEN.}
\noindent Man sucht wieder die Richtung, in der man das neue Werthen-System ändern muss, um \(R^2\) am schnellsten abnehmen zu lassen, ändert dasselbe dann demgemäss, soweit es die Ungleichungen zulassen, etc. und wiederholt das Verfahren so oft, bis endlich
\[
\delta(R^2)\geqq0
\]
\noindent wird, in welchem Falle die wirklichen Minimalwerthe der Grössen \(x_1,.\ .\ .,x_n\) gefunden sind.
